\begin{abstract}
La convergencia entre conectividad Direct-to-Device (D2D) y Redes No Terrestres (NTN) representa un pilar fundamental para el ecosistema 6G, habilitando cobertura ubicua, resiliencia y nuevos servicios para dispositivos estándar. Este artículo presenta un marco técnico integral que analiza arquitecturas híbridas TN-NTN, evoluciones de estándares 3GPP (Rel-17 a Rel-19 y perspectivas 6G), y desafíos técnicos clave como compensación Doppler, movilidad, gestión de interferencias y optimización de recursos. Se propone una arquitectura de referencia unificada con payloads regenerativos, orquestación inteligente TN/NTN y soporte nativo para D2D en bandas sub-7 GHz y Ka. Se discuten métricas de desempeño, simulaciones de cobertura y trade-offs entre throughput, latencia y consumo energético. Los resultados destacan mejoras significativas en disponibilidad de servicio y eficiencia espectral, estableciendo una base reproducible para el despliegue de redes 6G verdaderamente integradas y resilientes.
\end{abstract}